% ============================================================
% APPENDIX A: KEY SOURCE CODE MODULES
% ============================================================

\chapter{KEY SOURCE CODE MODULES}
\label{app:source_code}

This appendix presents the complete source code of the most critical modules in the QoS Scheduler system. These modules form the core of the Data Plane pipeline and are referenced throughout Chapters~4 and~5.

% -----------------------------------------------------------

\section{Packet Parser --- \texttt{RawPacket.kt}}
\label{app:rawpacket}

The \texttt{RawPacket} module is responsible for parsing raw byte arrays read from the TUN interface into structured packet objects. It supports both IPv4 and IPv6 with full Extension Header Chaining.

\begin{lstlisting}[language=Java, caption={RawPacket.kt --- Complete packet parser}, label={lst:rawpacket}]
package com.qos.scheduler.model

import java.nio.ByteBuffer

data class RawPacket(
    val ipVersion: Int,
    val srcIp: String, val dstIp: String,
    val srcPort: Int, val dstPort: Int,
    val protocol: Protocol,
    val length: Int,
    val ipHeaderLength: Int,
    val transportHeaderLength: Int,
    val payloadOffset: Int,
    val payloadLength: Int,
    val tcpSequenceNumber: Long? = null,
    val tcpAckNumber: Long? = null,
    val tcpWindowSize: Int? = null,
    val tcpFlags: TcpControlFlags? = null,
    val rawBuffer: ByteArray
) {
    val isIpv4: Boolean get() = ipVersion == 4
    val payload: ByteArray
        get() = rawBuffer.copyOfRange(
            payloadOffset, payloadOffset + payloadLength)

    companion object {
        fun parse(buffer: ByteArray, length: Int): RawPacket? {
            if (length < 20) return null
            val buf = ByteBuffer.wrap(buffer, 0, length)
            val version = (buf.get(0).toInt() and 0xFF) shr 4
            return when (version) {
                4 -> parseIpv4(buf, buffer, length)
                6 -> parseIpv6(buf, buffer, length)
                else -> null
            }
        }

        private fun parseIpv4(
            buf: ByteBuffer, raw: ByteArray, length: Int
        ): RawPacket? {
            val ihl = (buf.get(0).toInt() and 0x0F) * 4
            if (length < ihl + 4) return null
            val protocolNum = buf.get(9).toInt() and 0xFF
            val protocol = Protocol.fromNumber(protocolNum)
            val srcIp = buildString {
                for (i in 12..15) {
                    if (i > 12) append('.')
                    append(buf.get(i).toInt() and 0xFF)
                }
            }
            val dstIp = buildString {
                for (i in 16..19) {
                    if (i > 16) append('.')
                    append(buf.get(i).toInt() and 0xFF)
                }
            }
            val (srcPort, dstPort) = extractPorts(
                buf, ihl, protocol)
            val tcpMeta = extractTcpMetadata(
                buf, ihl, protocol, length)
            val thl = when (protocol) {
                Protocol.TCP -> tcpMeta.headerLength
                Protocol.UDP -> 8
                else -> 0
            }
            val po = (ihl + thl).coerceAtMost(length)
            val pl = (length - po).coerceAtLeast(0)
            return RawPacket(
                ipVersion = 4, srcIp = srcIp,
                dstIp = dstIp, srcPort = srcPort,
                dstPort = dstPort, protocol = protocol,
                length = length, ipHeaderLength = ihl,
                transportHeaderLength = thl,
                payloadOffset = po, payloadLength = pl,
                tcpSequenceNumber = tcpMeta.sequenceNumber,
                tcpAckNumber = tcpMeta.ackNumber,
                tcpWindowSize = tcpMeta.windowSize,
                tcpFlags = tcpMeta.flags,
                rawBuffer = raw.copyOf(length)
            )
        }

        private fun parseIpv6(
            buf: ByteBuffer, raw: ByteArray, length: Int
        ): RawPacket? {
            if (length < 40) return null
            var nextHeader = buf.get(6).toInt() and 0xFF
            var headerOffset = 40
            while (isIpv6ExtensionHeader(nextHeader)
                && headerOffset < length) {
                if (headerOffset + 2 > length) return null
                val extNext =
                    buf.get(headerOffset).toInt() and 0xFF
                val extLength = when (nextHeader) {
                    44 -> 8 // Fragment: fixed 8 bytes
                    else -> (buf.get(headerOffset + 1)
                        .toInt() and 0xFF) * 8 + 8
                }
                nextHeader = extNext
                headerOffset += extLength
            }
            val protocol = Protocol.fromNumber(nextHeader)
            val srcIp = formatIpv6(buf, 8)
            val dstIp = formatIpv6(buf, 24)
            val (srcPort, dstPort) = extractPorts(
                buf, headerOffset, protocol)
            val thl = if (protocol == Protocol.TCP
                && headerOffset + 13 < length) {
                ((buf.get(headerOffset + 12).toInt()
                    ushr 4) and 0x0F) * 4
            } else 8
            val po = (headerOffset+thl).coerceAtMost(length)
            val pl = (length - po).coerceAtLeast(0)
            return RawPacket(
                ipVersion = 6, srcIp = srcIp,
                dstIp = dstIp, srcPort = srcPort,
                dstPort = dstPort, protocol = protocol,
                length = length,
                ipHeaderLength = headerOffset,
                transportHeaderLength = thl,
                payloadOffset = po, payloadLength = pl,
                rawBuffer = raw.copyOf(length)
            )
        }

        private fun isIpv6ExtensionHeader(nh: Int) =
            nh in setOf(0, 43, 44, 60)

        private fun extractPorts(
            buf: ByteBuffer, off: Int, proto: Protocol
        ): Pair<Int, Int> {
            if (proto == Protocol.OTHER) return Pair(0, 0)
            if (buf.capacity() < off + 4) return Pair(0, 0)
            val src = ((buf.get(off).toInt() and 0xFF) shl 8
                ) or (buf.get(off+1).toInt() and 0xFF)
            val dst = ((buf.get(off+2).toInt() and 0xFF
                ) shl 8) or (buf.get(off+3).toInt() and 0xFF)
            return Pair(src, dst)
        }

        private fun extractTcpMetadata(
            buf: ByteBuffer, off: Int,
            proto: Protocol, len: Int
        ): TcpMetadata {
            if (proto != Protocol.TCP || off+20 > len)
                return TcpMetadata(0, null, null, null, null)
            val seq = buf.getInt(off+4).toLong() and
                0xFFFFFFFFL
            val ack = buf.getInt(off+8).toLong() and
                0xFFFFFFFFL
            val dataOff = ((buf.get(off+12).toInt()
                ushr 4) and 0x0F) * 4
            val flagsByte = buf.get(off+13).toInt() and 0xFF
            val win = buf.getShort(off+14).toInt() and 0xFFFF
            val flags = TcpControlFlags(
                fin = (flagsByte and 0x01) != 0,
                syn = (flagsByte and 0x02) != 0,
                rst = (flagsByte and 0x04) != 0,
                psh = (flagsByte and 0x08) != 0,
                ack = (flagsByte and 0x10) != 0
            )
            return TcpMetadata(dataOff, seq, ack, win, flags)
        }

        // IPv6 address formatting with :: compression
        private fun formatIpv6(
            buf: ByteBuffer, offset: Int
        ): String {
            val segs = IntArray(8) { i ->
                val h = buf.get(offset+i*2).toInt() and 0xFF
                val l = buf.get(offset+i*2+1).toInt() and 0xFF
                (h shl 8) or l
            }
            var maxStart = -1; var maxLen = 0
            var curStart = -1; var curLen = 0
            for (i in segs.indices) {
                if (segs[i] == 0) {
                    if (curStart == -1) {
                        curStart = i; curLen = 1
                    } else curLen++
                } else {
                    if (curLen > maxLen) {
                        maxStart = curStart
                        maxLen = curLen
                    }
                    curStart = -1; curLen = 0
                }
            }
            if (curLen > maxLen) {
                maxStart = curStart; maxLen = curLen
            }
            return buildString {
                var i = 0
                var justEmitted = false
                while (i < 8) {
                    if (i == maxStart && maxLen > 1) {
                        append("::"); i += maxLen
                        justEmitted = true
                        if (i >= 8) break
                    } else {
                        if (i > 0 && !justEmitted) append(':')
                        justEmitted = false
                        append(String.format("%x", segs[i]))
                        i++
                    }
                }
            }
        }

        private data class TcpMetadata(
            val headerLength: Int,
            val sequenceNumber: Long?,
            val ackNumber: Long?,
            val windowSize: Int?,
            val flags: TcpControlFlags?
        )
    }
}

data class TcpControlFlags(
    val fin: Boolean = false,
    val syn: Boolean = false,
    val rst: Boolean = false,
    val psh: Boolean = false,
    val ack: Boolean = false
)
\end{lstlisting}

% -----------------------------------------------------------

\section{Token Bucket --- \texttt{TokenBucket.kt}}
\label{app:tokenbucket}

The Token Bucket module implements the rate-limiting algorithm with nanosecond-precision timing. Thread safety is ensured via the \texttt{@Synchronized} annotation.

\begin{lstlisting}[language=Java, caption={TokenBucket.kt --- Complete rate limiter}, label={lst:tokenbucket}]
package com.qos.scheduler.scheduler

class TokenBucket(
    @Volatile var rateBps: Long,
    @Volatile var burstBytes: Long
) {
    private var tokens: Double = burstBytes.toDouble()
    private var lastRefillTime: Long = System.nanoTime()

    @Synchronized
    fun consume(bytes: Int): Boolean {
        refill()
        return if (tokens >= bytes) {
            tokens -= bytes
            true
        } else {
            false
        }
    }

    @Synchronized
    fun setRate(newRateBps: Long) {
        rateBps = newRateBps
        burstBytes = newRateBps * 2
        tokens = minOf(tokens, burstBytes.toDouble())
    }

    private fun refill() {
        val now = System.nanoTime()
        val elapsed =
            (now - lastRefillTime) / 1_000_000_000.0
        tokens = minOf(burstBytes.toDouble(),
                       tokens + rateBps * elapsed)
        lastRefillTime = now
    }
}
\end{lstlisting}

% -----------------------------------------------------------

\section{Bandwidth Scheduler --- \texttt{BandwidthScheduler.kt}}
\label{app:scheduler}

The Bandwidth Scheduler orchestrates all Token Buckets and implements the Weighted Fair Queuing rebalancing algorithm.

\begin{lstlisting}[language=Java, caption={BandwidthScheduler.kt --- WFQ orchestrator}, label={lst:scheduler}]
package com.qos.scheduler.scheduler

import com.qos.scheduler.model.AppTraffic
import com.qos.scheduler.model.TrafficClass
import java.util.concurrent.ConcurrentHashMap

class BandwidthScheduler {
    var uplinkBps: Long = 50_000_000L
    private val buckets =
        ConcurrentHashMap<String, TokenBucket>()
    private val manualCaps =
        ConcurrentHashMap<String, Long>()

    companion object {
        fun appBucketKey(uid: Int) = "app_$uid"
    }

    fun processPacket(key: String, size: Int): Boolean {
        val bucket = buckets[key] ?: return true
        return bucket.consume(size)
    }

    fun ensureBucket(
        key: String, priority: TrafficClass = TrafficClass.MEDIUM
    ) {
        if (!buckets.containsKey(key)) {
            val (rate, burst) = rateForClass(priority)
            buckets[key] = TokenBucket(rate, burst)
        }
    }

    fun setManualCap(ip: String, rateBps: Long) {
        manualCaps[ip] = rateBps
        buckets[ip]?.setRate(rateBps)
    }

    fun rebalanceWithApps(apps: Collection<AppTraffic>) {
        val hc = apps.count {
            it.priorityClass == TrafficClass.HIGH }
        val mc = apps.count {
            it.priorityClass == TrafficClass.MEDIUM }
        val lc = apps.count {
            it.priorityClass == TrafficClass.LOW }
        val totalWeight = hc*4 + mc*2 + lc*1 + 2
        if (totalWeight <= 0) return

        apps.forEach { app ->
            val key = appBucketKey(app.uid)
            if (manualCaps.containsKey(key)) return@forEach
            val weight = when (app.priorityClass) {
                TrafficClass.HIGH -> 4
                TrafficClass.MEDIUM -> 2
                TrafficClass.LOW -> 1
            }
            val rate = (uplinkBps * weight) / totalWeight
            buckets.getOrPut(key) {
                TokenBucket(rate, rate * 5) }.setRate(rate)
        }

        val hostRate = (uplinkBps * 2) / totalWeight
        buckets.getOrPut("__host__") {
            TokenBucket(hostRate, hostRate * 5)
        }.setRate(hostRate)
    }

    private fun rateForClass(
        cls: TrafficClass
    ): Pair<Long, Long> {
        val rate = when (cls) {
            TrafficClass.HIGH -> (uplinkBps * 0.8).toLong()
            TrafficClass.MEDIUM -> (uplinkBps * 0.5).toLong()
            TrafficClass.LOW -> (uplinkBps * 0.2).toLong()
        }
        val burst = when (cls) {
            TrafficClass.HIGH -> rate * 2
            TrafficClass.MEDIUM -> rate
            TrafficClass.LOW -> rate / 2
        }
        return Pair(rate, burst)
    }
}
\end{lstlisting}

% -----------------------------------------------------------

\section{DPI Classifier --- \texttt{DpiClassifier.kt}}
\label{app:dpi}

The DPI-Lite classifier categorizes traffic based on destination port heuristics.

\begin{lstlisting}[language=Java, caption={DpiClassifier.kt --- Port-based traffic classifier}, label={lst:dpi}]
package com.qos.scheduler.classifier

import com.qos.scheduler.model.Protocol
import com.qos.scheduler.model.RawPacket
import com.qos.scheduler.model.TrafficCategory

class DpiClassifier {
    fun classify(packet: RawPacket): TrafficCategory {
        val port = packet.dstPort
        val protocol = packet.protocol

        // VoIP / Real-time (High Priority)
        if (protocol == Protocol.UDP) {
            if (port in 5060..5061
                || port in 10000..20000)
                return TrafficCategory.VOIP
        }

        // Gaming (High Priority)
        val gamingPorts = listOf(
            27015, 27016, 5000, 5500, 8001, 8002)
        if (gamingPorts.contains(port))
            return TrafficCategory.ONLINE_GAMING

        // DNS
        if (port == 53)
            return TrafficCategory.WEB_BROWSING

        // Video Streaming
        if (port == 1935 || port == 8080)
            return TrafficCategory.STREAMING

        // Web / Social
        if (port == 80 || port == 443)
            return TrafficCategory.WEB_BROWSING

        // Downloads (Low Priority)
        val dlPorts = listOf(21, 22, 143, 993, 110, 995)
        if (dlPorts.contains(port))
            return TrafficCategory.FILE_TRANSFER

        return TrafficCategory.UNKNOWN
    }
}
\end{lstlisting}

% -----------------------------------------------------------

\section{Data Plane Processor --- \texttt{DataPlaneProcessor.kt}}
\label{app:dataplane}

The Data Plane Processor orchestrates the entire per-packet pipeline: parsing, UID resolution, classification, and scheduling decisions.

\begin{lstlisting}[language=Java, caption={DataPlaneProcessor.kt --- Pipeline orchestrator}, label={lst:dataplane}]
package com.qos.scheduler.service.dataplane

import com.qos.scheduler.classifier.DpiClassifier
import com.qos.scheduler.model.*
import com.qos.scheduler.scheduler.BandwidthScheduler
import com.qos.scheduler.util.AppResolver
import java.util.concurrent.ConcurrentHashMap

class DataPlaneProcessor(
    private val classifier: DpiClassifier,
    private val scheduler: BandwidthScheduler,
    private val appResolver: AppResolver,
    private val appTraffic:
        ConcurrentHashMap<Int, AppTraffic>,
    private val hostBucketKey: String
) {
    companion object {
        private val flowCache =
            ConcurrentHashMap<FlowKey, Int>()
    }

    data class Result(
        val parsed: Boolean,
        val droppedByScheduler: Boolean,
        val qosKey: String,
        val packet: RawPacket? = null
    )

    fun process(
        bytes: ByteArray, length: Int, isOutbound: Boolean
    ): Result {
        val pkt = RawPacket.parse(bytes, length)
            ?: return Result(false, false, hostBucketKey)

        val flowKey = if (isOutbound)
            FlowKey(pkt.srcIp, pkt.dstIp,
                pkt.srcPort, pkt.dstPort, pkt.protocol)
        else
            FlowKey(pkt.dstIp, pkt.srcIp,
                pkt.dstPort, pkt.srcPort, pkt.protocol)

        // STEP 1: UID Resolution (~1ms syscall)
        var uid = flowCache[flowKey]
        if (uid == null) {
            val pNum = if (pkt.protocol == Protocol.TCP)
                6 else 17
            val resolved = if (isOutbound)
                appResolver.getUidForConnection(pNum,
                    pkt.srcIp, pkt.srcPort,
                    pkt.dstIp, pkt.dstPort)
            else appResolver.getUidForConnection(pNum,
                    pkt.dstIp, pkt.dstPort,
                    pkt.srcIp, pkt.srcPort)

            if (resolved != null && resolved > 0) {
                flowCache[flowKey] = resolved
                uid = resolved
                // Retroactive migration of synthetic
                // port buckets occurs here (omitted
                // for brevity)
            }
        }

        // STEP 2: Classify & Track
        val cat = classifier.classify(pkt)
        var sKey = hostBucketKey
        if (uid != null && uid > 0) {
            val app = appTraffic.getOrPut(uid) {
                val info = appResolver.getAppInfo(uid)
                AppTraffic(uid, info.appName,
                    info.packageName, TrafficClass.MEDIUM)
            }
            sKey = BandwidthScheduler.appBucketKey(uid)
            if (isOutbound) app.bytesOut += length
            else app.bytesIn += length
            scheduler.ensureBucket(sKey, app.priorityClass)
        }

        // STEP 3: Scheduler Decision
        val allowed = if (isOutbound)
            scheduler.processPacket(sKey, length) else true

        return Result(true, !allowed, sKey, pkt)
    }
}
\end{lstlisting}
